% answers.tex —— 1992 年·三南卷 · 参考答案与详细解析片段
% 由 tools/generate_projects.py 依来源总账自 1992/1992_数学真题_三南_试题.md + 1992/1992_数学真题_三南_解析正文.md + 1992/1992_数学真题_三南_答案要点.md 生成，勿手改（改源 MD 后重跑）


\begin{gknotice}
\textbf{解析归属}：固定源 PDF 只有题面，没有答案页。以下推导依据固定题面形成，\textbf{非官方解析}；已由独立解题者在只读题面的条件下逐题盲解，结论与本节一致（第 15 题按命题意图取 D，题面缺损见试题卷说明）。原页疑点保留在题面与核源记录中，不以推导结果反改原题。
\end{gknotice}

\begin{answer}
\gksolhead{1}{}

\textbf{题目}：

设函数 $z=\mathrm{i}^2+\sqrt3\,\mathrm{i}$，那么 $\arg z$ 是\mbox{（　　）}。

（A）$\dfrac{5\pi}{6}$　　（B）$\dfrac{\pi}{3}$　　（C）$\dfrac{2\pi}{3}$　　（D）$-\dfrac{4\pi}{3}$

\Needspace{4\baselineskip}
\textbf{解析}

$z=\mathrm{i}^2+\sqrt3\,\mathrm{i}=-1+\sqrt3\,\mathrm{i}$ 对应第二象限的点，且 $|z|=2$，故主辐角 $\arg z=\pi-\dfrac{\pi}{3}=\dfrac{2\pi}{3}$，选 C。原卷用小写 $\arg$ 而未声明主值区间；若按「全体辐角」理解，D 的 $-\dfrac{4\pi}{3}$ 与 $\dfrac{2\pi}{3}$ 终边相同但不是主值，本题按通行约定取主值作答。

\gkanswer{答案：C（按辐角主值约定；见下注）}

\end{answer}
\begin{answer}
\gksolhead{2}{}

\textbf{题目}：

如果等边圆柱（即底面直径与母线相等的圆柱）的体积是 $16\pi\,\mathrm{cm}^3$，那么它的底面半径等于\mbox{（　　）}。

（A）$4\sqrt[3]{2}\,\mathrm{cm}$　　（B）$4\,\mathrm{cm}$　　（C）$2\sqrt[3]{2}\,\mathrm{cm}$　　（D）$2\,\mathrm{cm}$

\Needspace{4\baselineskip}
\textbf{解析}

等边圆柱的母线等于底面直径，设底面半径为 $r$，则高 $h=2r$，体积 $V=\pi r^2\cdot2r=2\pi r^3=16\pi$，故 $r^3=8$，$r=2\,\mathrm{cm}$，选 D。

\gkanswer{答案：D}

\end{answer}
\begin{answer}
\gksolhead{3}{}

\textbf{题目}：

$\dfrac{\arcsin\dfrac{\sqrt3}{2}-\arccos\left(-\dfrac12\right)}{\arctan(-\sqrt3)}$ 的值等于\mbox{（　　）}。

（A）$1$　　（B）$0$　　（C）$-\dfrac25$　　（D）$-\dfrac65$

\Needspace{4\baselineskip}
\textbf{解析}

$\arcsin\dfrac{\sqrt3}{2}=\dfrac{\pi}{3}$，$\arccos\left(-\dfrac12\right)=\dfrac{2\pi}{3}$，$\arctan(-\sqrt3)=-\dfrac{\pi}{3}$。分子为 $\dfrac{\pi}{3}-\dfrac{2\pi}{3}=-\dfrac{\pi}{3}$，与分母相同，比值为 $1$，选 A。

\gkanswer{答案：A}

\end{answer}
\begin{answer}
\gksolhead{4}{}

\textbf{题目}：

函数 $y=\log_{1/2}(1-x)\ (x<1)$ 的反函数是\mbox{（　　）}。

（A）$y=1+2^{-x}\ (x\in\mathbb R)$　　（B）$y=1-2^{-x}\ (x\in\mathbb R)$

（C）$y=1+2^x\ (x\in\mathbb R)$　　（D）$y=1-2^x\ (x\in\mathbb R)$

\Needspace{4\baselineskip}
\textbf{解析}

由 $y=\log_{1/2}(1-x)$ 得 $\left(\dfrac12\right)^y=1-x$，即 $x=1-2^{-y}$。交换 $x$、$y$ 得反函数 $y=1-2^{-x}$；原函数的值域为 $\mathbb R$，故反函数的定义域为 $\mathbb R$，选 B。

\gkanswer{答案：B}

\end{answer}
\begin{answer}
\gksolhead{5}{}

\textbf{题目}：

在长方体 $ABCD-A_1B_1C_1D_1$ 中，如果 $AB=BC=a$，$AA_1=2a$，那么点 $A$ 到直线 $A_1C$ 的距离等于\mbox{（　　）}。

（A）$\dfrac{2\sqrt6}{3}a$　　（B）$\dfrac{3\sqrt6}{2}a$　　（C）$\dfrac{2\sqrt3}{3}a$　　（D）$\dfrac{\sqrt6}{3}a$

\Needspace{4\baselineskip}
\textbf{解析}

\Needspace{5\baselineskip}
$AA_1\perp$ 底面，故 $AA_1\perp AC$，$\triangle AA_1C$ 在 $A$ 处为直角，两直角边 $AA_1=2a$、$AC=\sqrt2\,a$，斜边 $A_1C=\sqrt6\,a$。点 $A$ 到斜边 $A_1C$ 的距离即斜边上的高
\nopagebreak[4]
\[
d=\frac{AA_1\cdot AC}{A_1C}=\frac{2a\cdot\sqrt2\,a}{\sqrt6\,a}=\frac{2\sqrt3}{3}a,
\]
选 C。

\gkanswer{答案：C}

\end{answer}
\begin{answer}
\gksolhead{6}{}

\textbf{题目}：

函数 $y=\sin x\cos x+\sqrt3\cos^2x-\dfrac{\sqrt3}{2}$ 的最小正周期等于\mbox{（　　）}。

（A）$\pi$　　（B）$2\pi$　　（C）$\dfrac\pi4$　　（D）$\dfrac\pi2$

\Needspace{4\baselineskip}
\textbf{解析}

\nopagebreak[4]
\[
\begin{aligned}
y&=\frac12\sin2x+\sqrt3\cdot\frac{1+\cos2x}{2}-\frac{\sqrt3}{2} \\
&=\frac12\sin2x+\frac{\sqrt3}{2}\cos2x=\sin\left(2x+\frac{\pi}{3}\right),
\end{aligned}
\]
最小正周期 $T=\dfrac{2\pi}{2}=\pi$，选 A。

\gkanswer{答案：A}

\end{answer}
\begin{answer}
\gksolhead{7}{}

\textbf{题目}：

有一个椭圆，它的极坐标方程是\mbox{（　　）}。

（A）$\rho=\dfrac5{\sqrt3-2\cos\theta}$　　（B）$\rho=\dfrac5{\sqrt3-\sqrt3\cos\theta}$

（C）$\rho=\dfrac{2-\sqrt3\cos\theta}{5}$　　（D）$\rho=\dfrac5{2-\sqrt3\cos\theta}$

\Needspace{4\baselineskip}
\textbf{解析}

以焦点为极点的圆锥曲线方程为 $\rho=\dfrac{l}{1-e\cos\theta}$。选项 D 化为 $\rho=\dfrac{5/2}{1-\frac{\sqrt3}{2}\cos\theta}$，离心率 $e=\dfrac{\sqrt3}{2}<1$，是椭圆；A 的离心率 $\dfrac{2}{\sqrt3}>1$ 为双曲线，B 的离心率为 $1$ 是抛物线，C 不符该形式，故选 D。

\gkanswer{答案：D}

\end{answer}
\begin{answer}
\gksolhead{8}{}

\textbf{题目}：

不等式 $\left|\sqrt{x-2}-3\right|<1$ 的解集是\mbox{（　　）}。

（A）$\{x\mid 5<x<16\}$　　（B）$\{x\mid 6<x<18\}$　　（C）$\{x\mid 7<x<20\}$　　（D）$\{x\mid 8<x<22\}$

\Needspace{4\baselineskip}
\textbf{解析}

$\left|\sqrt{x-2}-3\right|<1$ 等价于 $2<\sqrt{x-2}<4$，平方得 $4<x-2<16$，即 $6<x<18$，选 B。

\gkanswer{答案：B}

\end{answer}
\begin{answer}
\gksolhead{9}{}

\textbf{题目}：

设等差数列 $\{a_n\}$ 的公差是 $d$，如果它的前 $n$ 项和 $S_n=-n^2$，那么\mbox{（　　）}。

（A）$a_n=2n-1,\ d=-2$　　（B）$a_n=2n-1,\ d=2$

（C）$a_n=-2n+1,\ d=-2$　　（D）$a_n=-2n+1,\ d=2$

\Needspace{4\baselineskip}
\textbf{解析}

$a_1=S_1=-1$；当 $n\ge2$ 时 $a_n=S_n-S_{n-1}=-n^2+(n-1)^2=-2n+1$，该式对 $n=1$ 也成立，故 $a_n=-2n+1$，公差 $d=a_2-a_1=-2$，选 C。

\gkanswer{答案：C}

\end{answer}
\begin{answer}
\gksolhead{10}{}

\textbf{题目}：

方程 $\cos2x=3\cos x+1$ 的解集是\mbox{（　　）}。

（A）$\{x\mid x=2k\pi\pm\dfrac{2\pi}{3},\ k\in\mathbb Z\}$　　（B）$\{x\mid x=k\pi\pm\dfrac\pi3,\ k\in\mathbb Z\}$

（C）$\{x\mid x=k\pi\pm\dfrac{2\pi}{3},\ k\in\mathbb Z\}$　　（D）$\{x\mid x=2k\pi\pm\dfrac\pi3,\ k\in\mathbb Z\}$

\Needspace{4\baselineskip}
\textbf{解析}

令 $c=\cos x$，得 $2c^2-1=3c+1$，即 $2c^2-3c-2=0$，$(2c+1)(c-2)=0$。因 $c\in[-1,1]$，只能 $c=-\dfrac12$，故 $x=2k\pi\pm\dfrac{2\pi}{3}$（$k\in\mathbb Z$），选 A。

\gkanswer{答案：A}

\end{answer}
\begin{answer}
\gksolhead{11}{}

\textbf{题目}：

有一条半径是 $2$ 的弧，其度数是 $60^\circ$，它绕经过弧的中点的直径旋转得到一个球冠，那么这个球冠的面积是\mbox{（　　）}。

（A）$4(2-\sqrt3)\pi$　　（B）$2(2-\sqrt3\pi)$　　（C）$4\sqrt3\pi$　　（D）$2\sqrt3\pi$

\Needspace{4\baselineskip}
\textbf{解析}

半径 $R=2$ 的弧对应圆心角 $60^\circ$，绕经过弧中点的直径旋转所得球冠的角半径为 $30^\circ$，冠高 $h=R-R\cos30^\circ=2-\sqrt3$，面积 $S=2\pi Rh=4(2-\sqrt3)\pi$，选 A。原卷选项 B 印作 \(2(2-\sqrt3\pi)\)，$\pi$ 落在括号内，与其余选项量纲不一致，应为排印移位；按字面亦不为正确答案，故不影响选项判定。

\gkanswer{答案：A}

\end{answer}
\begin{answer}
\gksolhead{12}{}

\textbf{题目}：

某小组共有 $10$ 名学生，其中女生 $3$ 名。现选举 $2$ 名代表，至少有 $1$ 名女生当选的不同的选法共有\mbox{（　　）}。

（A）$27$ 种　　（B）$48$ 种　　（C）$21$ 种　　（D）$24$ 种

\Needspace{4\baselineskip}
\textbf{解析}

至少有 1 名女生的选法为 $\mathrm{C}_{10}^{2}-\mathrm{C}_{7}^{2}=45-21=24$，选 D。

\gkanswer{答案：D}

\end{answer}
\begin{answer}
\gksolhead{13}{}

\textbf{题目}：

设全集 $U=\mathbb R$，集合 $M=\{x\mid\sqrt{x^2}>2\}$，$N=\{x\mid\log_x7>\log_3 7\}$，那么 $M\cap\overline N=$\mbox{（　　）}。

（A）$\{x\mid x<-2\}$　　（B）$\{x\mid x<-2\text{ \text{或} }x\ge3\}$

（C）$\{x\mid x\ge3\}$　　（D）$\{x\mid -2\le x<3\}$

\Needspace{4\baselineskip}
\textbf{解析}

\Needspace{5\baselineskip}
$M=\{x\mid\sqrt{x^2}>2\}=\{x\mid|x|>2\}=(-\infty,-2)\cup(2,+\infty)$。对 $N$ 须 $x>0$ 且 $x\ne1$：当 $1<x<3$ 时 $\log_x7>\log_37$；当 $0<x<1$ 时 $\log_x7<0<\log_37$；当 $x>3$ 时结论反向。故 $N=(1,3)$，$\overline N=(-\infty,1]\cup[3,+\infty)$，
\nopagebreak[4]
\[
M\cap\overline N=(-\infty,-2)\cup[3,+\infty),
\]
选 B。

\gkanswer{答案：B}

\end{answer}
\begin{answer}
\gksolhead{14}{}

\textbf{题目}：

设 $\{a_n\}$ 是由正数组成的等比数列，公比 $q=2$，且 $a_1a_2a_3\cdots a_{30}=2^{30}$，那么 $a_3a_6a_9\cdots a_{30}$ 等于\mbox{（　　）}。

（A）$2^{10}$　　（B）$2^{20}$　　（C）$2^{16}$　　（D）$2^{15}$

\Needspace{4\baselineskip}
\textbf{解析}

\Needspace{5\baselineskip}
$a_n=a_1\cdot2^{n-1}$，故 $a_1a_2\cdots a_{30}=a_1^{30}\cdot2^{0+1+\cdots+29}=a_1^{30}\cdot2^{435}=2^{30}$，得 $a_1^{10}=2^{-135}$。又
\nopagebreak[4]
\[
a_3a_6\cdots a_{30}=a_1^{10}\cdot2^{2+5+\cdots+29}=2^{-135}\cdot2^{155}=2^{20},
\]
选 B。

\gkanswer{答案：B}

\end{answer}
\begin{answer}
\gksolhead{15}{}

\textbf{题目}：

设 $\triangle ABC$ 不是直角三角形，$A$ 和 $B$ 是它的两个内角，那么\mbox{（　　）}。

（A）“$A<B$”是“$\tan A<\tan B$”的充分条件，但不是必要条件

（B）“$A<B$”是“$\tan A<\tan B$”的必要条件，但不是充分条件

（C）“$A<B$”是“$\tan A<\tan B$”的充分必要条件

（D）“$A<B$”是“$\tan A<\tan B$”的充分条件，也不是必要条件

\Needspace{4\baselineskip}
\textbf{解析}

原卷四个选项中没有正确项：A、D 均声称「$A<B$」是「$\tan A<\tan B$」的充分条件，但取 $(A,B,C)=(50^\circ,120^\circ,10^\circ)$（非直角三角形）时 $A<B$ 而 $\tan A>0>\tan B$，不充分；又取 $(A,B,C)=(120^\circ,40^\circ,20^\circ)$ 时 $\tan A<0<\tan B$ 而 $A>B$，不必要。故该命题\textbf{既非充分也非必要}。D 项原印作「充分条件，也不是必要条件」，漏「既不」二字；按命题意图应补足为「既不是充分条件，也不是必要条件」，此时选 D。题面照录原样，不作静默补正。

\gkanswer{答案：D（条件结论：原卷 D 项漏印「既不」二字；按字面 A、D 均称充分而不必要，无正确选项，按命题意图补足后为 D）}

\end{answer}
\begin{answer}
\gksolhead{16}{}

\textbf{题目}：

对于定义域是 $\mathbb R$ 的任何奇函数 $f(x)$，都有\mbox{（　　）}。

（A）$f(x)-f(-x)>0\ (x\in\mathbb R)$　　（B）$f(x)-f(-x)\le0\ (x\in\mathbb R)$

（C）$f(x)f(-x)\le0\ (x\in\mathbb R)$　　（D）$f(x)f(-x)>0\ (x\in\mathbb R)$

\Needspace{4\baselineskip}
\textbf{解析}

$f$ 为奇函数，$f(-x)=-f(x)$，故 $f(x)f(-x)=-f(x)^2\le0$ 对一切 $x\in\mathbb R$ 成立（含 $f(x)=0$ 的情形），选 C。A、B 中的 $f(x)-f(-x)=2f(x)$ 符号不定，D 恒不成立。

\gkanswer{答案：C}

\end{answer}
\begin{answer}
\gksolhead{17}{}

\textbf{题目}：

如果双曲线的两条渐近线的方程是 $y=\pm\dfrac32x$，焦点坐标是 $(-\sqrt{26},0)$ 和 $(\sqrt{26},0)$，那么它的两条准线之间的距离是\mbox{（　　）}。

（A）$\dfrac8{13}\sqrt{26}$　　（B）$\dfrac4{13}\sqrt{26}$　　（C）$\dfrac{18}{13}\sqrt{26}$　　（D）$\dfrac9{13}\sqrt{26}$

\Needspace{4\baselineskip}
\textbf{解析}

设双曲线为 $\dfrac{x^2}{a^2}-\dfrac{y^2}{b^2}=1$，由渐近线斜率得 $\dfrac{b}{a}=\dfrac32$，由 $c=\sqrt{26}$ 得 $a^2+b^2=26$，即 $a^2\left(1+\dfrac94\right)=26$，$a^2=8$。两准线间距离 $\dfrac{2a^2}{c}=\dfrac{16}{\sqrt{26}}=\dfrac{8\sqrt{26}}{13}$，选 A。

\gkanswer{答案：A}

\end{answer}
\begin{answer}
\gksolhead{18}{}

\textbf{题目}：

$\tan\dfrac\pi8=\underline{\qquad\qquad}$。

\Needspace{4\baselineskip}
\textbf{解析}

设 $t=\tan\dfrac{\pi}{8}>0$，由 $\tan\dfrac{\pi}{4}=\dfrac{2t}{1-t^2}=1$ 得 $t^2+2t-1=0$，故 $t=\sqrt2-1$。

\gkanswer{答案：\(\sqrt2-1\)}

\end{answer}
\begin{answer}
\gksolhead{19}{}

\textbf{题目}：

设直线的参数方程是 $\begin{cases}x=2+\dfrac12t,\\y=3+\dfrac{\sqrt3}{2}t,\end{cases}$ 那么它的斜截式方程是\underline{\qquad\qquad}。

\Needspace{4\baselineskip}
\textbf{解析}

由 $x=2+\dfrac12t$ 得 $t=2(x-2)$，代入 $y=3+\dfrac{\sqrt3}{2}t$ 得 $y=3+\sqrt3(x-2)=\sqrt3x+3-2\sqrt3$。

\gkanswer{答案：\(y=\sqrt3x+3-2\sqrt3\)}

\end{answer}
\begin{answer}
\gksolhead{20}{}

\textbf{题目}：

如果三角形的顶点分别是 $O(0,0)$、$A(0,15)$、$B(-8,0)$，那么它的内切圆方程是\underline{\qquad\qquad}。

\Needspace{4\baselineskip}
\textbf{解析}

$\angle O=90^\circ$，三边为 $8$、$15$、$17$，内切圆半径 $r=\dfrac{8+15-17}{2}=3$。圆心在第二象限且与两坐标轴相切，故圆心为 $(-3,3)$，所求方程为 $(x+3)^2+(y-3)^2=9$。校验：到斜边 $15x-8y+120=0$ 的距离为 $\dfrac{|-45-24+120|}{17}=3$。

\gkanswer{答案：\((x+3)^2+(y-3)^2=9\)}

\end{answer}
\begin{answer}
\gksolhead{21}{}

\textbf{题目}：

$\displaystyle\lim_{n\to\infty}\left[\frac1{1\cdot4}+\frac1{4\cdot7}+\frac1{7\cdot10}+\cdots+\frac1{(3n-2)(3n+1)}\right]=$\underline{\qquad\qquad}。

\Needspace{4\baselineskip}
\textbf{解析}

\Needspace{5\baselineskip}
$\dfrac{1}{(3k-2)(3k+1)}=\dfrac13\left(\dfrac{1}{3k-2}-\dfrac{1}{3k+1}\right)$，故部分和
\nopagebreak[4]
\[
\frac13\left(1-\frac{1}{3n+1}\right)\to\frac13\quad(n\to\infty).
\]

\gkanswer{答案：\(\dfrac13\)}

\end{answer}
\begin{answer}
\gksolhead{22}{}

\textbf{题目}：

$91^{92}$ 除以 $100$ 的余数是\underline{\qquad\qquad}。

\Needspace{4\baselineskip}
\textbf{解析}

\Needspace{5\baselineskip}
$91^2\equiv81$，$91^4\equiv61$，$91^8\equiv21$，$91^{16}\equiv41$，$91^{32}\equiv81$，$91^{64}\equiv61\pmod{100}$。由 $92=64+16+8+4$，
\nopagebreak[4]
\[
91^{92}\equiv61\cdot41\cdot21\cdot61\equiv81\pmod{100},
\]
故余数为 $81$。

\gkanswer{答案：\(81\)}

\end{answer}
\begin{answer}
\gksolhead{23}{}

\textbf{题目}：

已知三棱锥 $A-BCD$ 的体积是 $V$，棱 $BC$ 的长是 $a$，面 $ABC$ 和面 $DBC$ 的面积分别是 $S_1$ 和 $S_2$。设面 $ABC$ 和面 $DBC$ 所成的二面角是 $\alpha$，那么 $\sin\alpha=$\underline{\qquad\qquad}。

\Needspace{4\baselineskip}
\textbf{解析}

\Needspace{5\baselineskip}
设 $A$、$D$ 到直线 $BC$ 的距离分别为 $h_1=\dfrac{2S_1}{a}$、$h_2=\dfrac{2S_2}{a}$，则二面角 $\alpha$ 即两垂线方向之夹角，三棱锥的高为 $h_2\sin\alpha$。以 $\triangle ABC$ 为底，
\nopagebreak[4]
\[
V=\frac13S_1h_2\sin\alpha=\frac13\cdot\frac{a h_1}{2}\cdot\frac{2S_2}{a}\sin\alpha=\frac{2S_1S_2}{3a}\sin\alpha,
\]
故 $\sin\alpha=\dfrac{3aV}{2S_1S_2}$。结果的 $\sin\alpha$ 对补角取同值，与二面角取锐角或钝角无关。

\gkanswer{答案：\(\dfrac{3aV}{2S_1S_2}\)}

\end{answer}
\begin{answer}
\gksolhead{24}{}

\textbf{题目}：

已知关于 $x$ 的方程 $2a^{2x-2}-7a^{x-1}+3=0$ 有一个根是 $2$，求 $a$ 的值和方程其余的根。

\Needspace{4\baselineskip}
\textbf{解析}

把 $x=2$ 代入得 $2a^2-7a+3=0$，即 $(2a-1)(a-3)=0$，故 $a=\dfrac12$ 或 $a=3$（均满足底数要求 $a>0$ 且 $a\ne1$）。

令 $t=a^{x-1}$，原方程化为 $2t^2-7t+3=0$，得 $t=\dfrac12$ 或 $t=3$。

\begin{itemize}
\item 若 $a=\dfrac12$：$\left(\dfrac12\right)^{x-1}=\dfrac12$ 得 $x=2$；$\left(\dfrac12\right)^{x-1}=3$ 得 $x=1-\log_23$。
\item 若 $a=3$：$3^{x-1}=3$ 得 $x=2$；$3^{x-1}=\dfrac12$ 得 $x=1-\log_32$。
\end{itemize}

故 $a=\dfrac12$ 时其余根为 $x=1-\log_23$；$a=3$ 时其余根为 $x=1-\log_32$。

\gkanswer{答案：\(a=\dfrac12\) 时另有根 \(x=1-\log_23\)；\(a=3\) 时另有根 \(x=1-\log_32\)。}

\end{answer}
\begin{answer}
\gksolhead{25}{}

\textbf{题目}：

已知平面 $\alpha$ 和不在这个平面内的直线 $a$ 都垂直于平面 $\beta$。求证：$a\parallel\alpha$。

\Needspace{4\baselineskip}
\textbf{解析}

由 $\alpha\perp\beta$ 知两平面相交，设交线为 $c$。在 $\alpha$ 内作 $b\perp c$，由面面垂直的性质定理得 $b\perp\beta$。又 $a\perp\beta$，垂直于同一平面的两条直线互相平行，故 $a\parallel b$。因 $b\subset\alpha$ 且题设 $a\not\subset\alpha$，由线面平行的判定定理得 $a\parallel\alpha$。

\gkanswer{答案：见解析（线面平行判定）。}

\end{answer}
\begin{answer}
\gksolhead{26}{}

\textbf{题目}：

证明不等式：$1+\dfrac1{\sqrt2}+\dfrac1{\sqrt3}+\cdots+\dfrac1{\sqrt n}<2\sqrt n\quad(n\in\mathbb N^*)$。

\Needspace{4\baselineskip}
\textbf{解析}

\Needspace{5\baselineskip}
对每个 $k\ge1$，
\nopagebreak[4]
\[
\frac{1}{\sqrt{k}}<\frac{2}{\sqrt{k}+\sqrt{k-1}}=2\left(\sqrt{k}-\sqrt{k-1}\right),
\]
\Needspace{5\baselineskip}
因为 $\sqrt{k}+\sqrt{k-1}<2\sqrt{k}$。逐项相加并利用右侧望远镜和，
\nopagebreak[4]
\[
1+\frac{1}{\sqrt2}+\frac{1}{\sqrt3}+\cdots+\frac{1}{\sqrt n}<2\sqrt n\quad(n\in\mathbb N^*).
\]

\gkanswer{答案：见解析（逐项放缩后望远镜求和）。}

\end{answer}
\begin{answer}
\gksolhead{27}{}

\textbf{题目}：

设抛物线经过两点 $(-1,6)$ 和 $(-1,-2)$ 对称轴与 $x$ 轴平行，开口向右，直线 $y=2x+7$ 被抛物线截得的线段的长是 $4\sqrt{10}$，求抛物线的方程。

\Needspace{4\baselineskip}
\textbf{解析}

设抛物线为 $(y-k)^2=2p(x-h)$，$p>0$。两点横坐标同为 $-1$，故 $(6-k)^2=(-2-k)^2$，得 $k=2$。代回得 $16=2p(-1-h)$，即 $h=-1-\dfrac{8}{p}$，于是 $(y-2)^2=2px+2p+16$。记 $q=2p>0$，则 $(y-2)^2=q(x+1)+16$。

\Needspace{5\baselineskip}
与 $y=2x+7$ 联立（$x=\dfrac{y-7}{2}$）消去 $x$：
\nopagebreak[4]
\[
2y^2-(8+q)y+(5q-24)=0.
\]
\Needspace{5\baselineskip}
设交点纵坐标为 $y_1$、$y_2$，则 $|\Delta y|^2=\dfrac{(8+q)^2-8(5q-24)}{4}$。弦长为
\nopagebreak[4]
\[
\sqrt{1+\frac14}\,|\Delta y|=\frac{\sqrt5}{2}|\Delta y|=4\sqrt{10},
\]
\Needspace{5\baselineskip}
故 $|\Delta y|^2=128$，得 $(8+q)^2-8(5q-24)=512$，即 $q^2-24q-256=0$，$q=32$（$q=-8$ 舍去），$p=16$。所求抛物线为
\nopagebreak[4]
\[
(y-2)^2=32x+48=32\left(x+\frac32\right),\quad\text{即}\quad y^2-4y-32x-44=0.
\]
校验：$(-1,6)$、$(-1,-2)$ 均在抛物线上；弦端纵坐标为 $10\pm4\sqrt2$，$|\Delta y|=8\sqrt2$，弦长 $\dfrac{\sqrt5}{2}\cdot8\sqrt2=4\sqrt{10}$，与题设相符。

\gkanswer{答案：\((y-2)^2=32\left(x+\dfrac32\right)\)，即 \(y^2-4y-32x-44=0\)。}

\end{answer}
\begin{answer}
\gksolhead{28}{}

\textbf{题目}：

求同时满足下列两个条件的所有复数 $z$：

① $z+\dfrac{10}{z}$ 是实数，且 $1<z+\dfrac{10}{z}\le6$；

② $z$ 的实部和虚部都是整数。

\Needspace{4\baselineskip}
\textbf{解析}

\Needspace{5\baselineskip}
设 $z=x+y\mathrm{i}$（$x,y\in\mathbb R$，$z\ne0$）。由
\nopagebreak[4]
\[
z+\frac{10}{z}=x+y\mathrm{i}+\frac{10(x-y\mathrm{i})}{x^2+y^2}
\]
的虚部为 $y\left(1-\dfrac{10}{x^2+y^2}\right)$，条件①为实数给出 $y=0$ 或 $x^2+y^2=10$。

若 $y=0$，则 $z=x\ne0$ 为实数，$z+\dfrac{10}{z}=x+\dfrac{10}{x}$：当 $x>0$ 时该值 $\ge2\sqrt{10}>6$，当 $x<0$ 时该值 $<0$，均不满足 $1<z+\dfrac{10}{z}\le6$。

\Needspace{5\baselineskip}
若 $x^2+y^2=10$，则 $\dfrac{10}{z}=\overline z$，$z+\dfrac{10}{z}=2x$，条件变为 $\dfrac12<x\le3$。由 $x,y$ 为整数且平方和为 $10$，得 $(x,y)=(\pm1,\pm3)$ 或 $(\pm3,\pm1)$；结合 $\dfrac12<x\le3$ 只需 $x=1$ 或 $x=3$，故
\nopagebreak[4]
\[
z=1\pm3\mathrm{i}\quad\text{或}\quad z=3\pm\mathrm{i}.
\]
校验：$|z|^2=10$，$z+\dfrac{10}{z}=2x$ 分别为 $2$ 与 $6$，均在 $(1,6]$ 内。

\gkanswer{答案：\(z=1\pm3\mathrm{i}\) 或 \(z=3\pm\mathrm{i}\)。}

\end{answer}